%! Author = Sujal Singh
%! Date = 12/19/23

% Preamble
\documentclass[11pt]{beamer}
\title{Equilibrium}
\author[Aman, Dhruv, Daksh...]
{\(|\) Aman Gupta \(|\) Dhruv Grover \(|\) Daksh Jain \(|\) Prashant Pulkit \(|\)\\
    \(|\) Pranav Bisht \(|\) Sujal Singh \(|\)}
\date[Prashant, Pranav, Sujal]{\textbf{Enrollment Numbers:}\\0\{00-00\}00000000}

\usetheme{Madrid}
\usecolortheme{beaver}
\setbeamertemplate{navigation symbols}{}
\setbeamertemplate{frametitle continuation}[from second][]
\setbeamertemplate{page number in head/foot}{}

% Packages
\usepackage{amsmath}
\usepackage{gensymb}
\usepackage[linewidth=0.75pt]{structmech}
\usepackage{tikz}

\usetikzlibrary{angles,arrows.meta}

\newcommand{\0}{\footnotesize}

% Document
\begin{document}
%    \begin{frame}{Mechanics Presentation (Group 2)}
%        \begin{center}
%            %! suppress = FileNotFound
%            \includegraphics[width=80pt]{logo}
%        \end{center}\vspace*{-10pt}
%        \maketitle
%    \end{frame}
%
%    \begin{frame}[t]{Force Body Diagram}
%        ~
%    \end{frame}
%
%    \begin{frame}[t]{Equations of Equilibrium}
%        ~
%    \end{frame}
%
%    \begin{frame}[t]{Applications to Engineering Problems}
%        ~
%    \end{frame}
%
%    \begin{frame}[t]{Equilibrium of Two Force Members}
%        ~
%    \end{frame}
%
%    \begin{frame}[t]{Equilibrium of Three Force Members}
%        ~
%    \end{frame}

    \begin{frame}[t,allowframebreaks]{Numericals}
        Two equal loads of $2500$ N are supported by a flexible string $ABCD$ at points $B$ and $C$.\ Find the
        tensions in the portions $AB$, $BC$ and $CD$ of the string.\\[20pt]

        \begin{center}
            \begin{tikzpicture}
                \FixedSupport[-90]{0,4}{1.5}
                \filldraw[black] (0,4) circle (1pt);
                \draw [thick] (0,4) coordinate (A) node[anchor=east,xshift=-5pt] {\0$A$}
                -- node[midway,anchor=south west,xshift=-4pt] {\0$T_1$} ++(-60:2) coordinate (B)
                node[anchor=north east] {\0$B$}
                -- ++(-30:2) node[midway,anchor=north east,xshift=5pt] {\0$T_2$} coordinate (C)
                node[anchor=north west] {\0$C$}
                -- ++(30:4) node[midway,anchor=north west] {\0$T_3$} coordinate (D)
                node[anchor=west,xshift=5pt] {\0$D$};
                \FixedSupport[90]{6.1961,3.2679}{1.5}
                \filldraw[black] (D) circle (1pt);
                \draw [dashed] (B) -- ++(0,1) coordinate (b);
                \draw [dashed] (C) -- ++(0,1) coordinate (c);
                \draw [-latex] (B) -- ++(0,-1) node[anchor=north] {\0 2500 N};
                \draw [-latex] (C) -- ++(0,-1) node[anchor=north] {\0 2500 N};
                \draw pic [draw,angle radius=8] {angle=C--B--b} node[yshift=2.5cm,xshift=1.6cm] {\tiny$120\degree$}
                \draw pic [draw,angle radius=8] {angle=c--C--B} node[yshift=1.75cm,xshift=2.5cm] {\tiny$60\degree$}
                \path (A) -- ++(0,-1) coordinate (a);
                \path (D) -- ++(0,-1) coordinate (d);
                \draw pic [draw,angle radius=8] {angle=C--D--d} node[yshift=2.8cm,xshift=5.9cm] {\tiny$60\degree$}
                \draw pic [draw,angle radius=8] {angle=a--A--B} node[yshift=3.3cm,xshift=0.1cm] {\tiny$30\degree$}
            \end{tikzpicture}
        \end{center}
    \end{frame}

%    \begin{frame}{End}
%        \begin{center}
%            \huge Thank you for listening.
%        \end{center}
%    \end{frame}
\end{document}
